\begin{notechapter}
\chapter{Lipschitz Continuity}\label{ch:lipschitz-continuity}
\chaptersubtitle{A speed limit for functions, and why machine learning needs one}

\begin{chapterintro}
Many guarantees in machine learning --- that a classifier cannot be fooled by
a tiny perturbation, that gradient descent with a given step size makes
progress, that a GAN critic estimates a meaningful distance --- rest on one
property: the function involved cannot change faster than a known rate. You
will learn
\begin{itemize}[nosep]
  \item the definition of Lipschitz continuity and its geometric meaning;
  \item how to compute or bound Lipschitz constants: gradients, compositions,
        linear layers and whole networks;
  \item how a Lipschitz bound certifies robustness to adversarial
        perturbations;
  \item the descent lemma: why a Lipschitz \emph{gradient} fixes the
        admissible learning rate;
  \item the role of 1-Lipschitz functions in Wasserstein GANs, and other uses.
\end{itemize}
\prerequisites{multivariable calculus (gradient, Jacobian, chain rule);
matrix norms and singular values.}
\end{chapterintro}

\section{Definition and intuition}\label{lipschitz-continuity:sec:definition}

\begin{definition}[Lipschitz continuity]\label{lipschitz-continuity:def:lipschitz}
A function $f:\R^n\to\R^m$ is \emph{$L$-Lipschitz} (for norms $\norm{\cdot}$ on
both spaces) if
\begin{equation}\label{lipschitz-continuity:eq:lipschitz}
  \norm{f(x_1) - f(x_2)} \leq L\,\norm{x_1 - x_2}\qquad\text{for all } x_1,x_2.
\end{equation}
The smallest such $L$ is the \emph{Lipschitz constant} of $f$, written
$\mathrm{Lip}(f)$. Unless stated otherwise we use Euclidean norms.
\end{definition}

\begin{intuition}[A speed limit, and a cone]
Dividing by $\norm{x_1-x_2}$, the definition says that no ``rate of change''
of $f$ --- no slope of a chord of its graph --- exceeds $L$. For a function of
one variable this has a picture: at every point of the graph, place a double
cone with slopes $\pm L$; the graph never enters the inside of the cone
(\cref{lipschitz-continuity:fig:cone}). A small $L$ means a function that
changes slowly; a large $L$ allows sharp variations.
\end{intuition}

\begin{figure}[ht]
\centering
\begin{tikzpicture}
\begin{axis}[width=0.62\textwidth, height=7cm, xmin=-3, xmax=3, ymin=-2.2, ymax=2.2,
  axis lines=middle, xtick=\empty, ytick=\empty, samples=120]
  \addplot[name path=up, draw=none, domain=-3:3] {0.5*sin(deg(2*0.6)) + 0.3*0.6 + 1.3*abs(x-0.6)};
  \addplot[name path=dn, draw=none, domain=-3:3] {0.5*sin(deg(2*0.6)) + 0.3*0.6 - 1.3*abs(x-0.6)};
  \addplot[name path=top, draw=none, domain=-3:3] {3};
  \addplot[name path=bot, draw=none, domain=-3:3] {-3};
  \addplot[bookpitfall!15] fill between[of=up and top];
  \addplot[bookpitfall!15] fill between[of=dn and bot];
  \addplot[bookpitfall, dashed, domain=-3:3] {0.5*sin(deg(2*0.6)) + 0.3*0.6 + 1.3*(x-0.6)};
  \addplot[bookpitfall, dashed, domain=-3:3] {0.5*sin(deg(2*0.6)) + 0.3*0.6 - 1.3*(x-0.6)};
  \addplot[very thick, bookaccent, domain=-3:3] {0.5*sin(deg(2*x)) + 0.3*x};
  \fill (axis cs:0.6,{0.5*sin(deg(1.2)) + 0.18}) circle (2pt);
\end{axis}
\end{tikzpicture}
\caption{The function $f(x) = \frac12\sin 2x + 0.3x$ has $\abs{f'}\leq 1.3$, so it
is $1.3$-Lipschitz: at every point (here the black dot) its graph stays
outside the shaded region above and below the lines of slope $\pm1.3$.}
\label{lipschitz-continuity:fig:cone}
\end{figure}

\begin{example}[Lipschitz or not?]\label{lipschitz-continuity:ex:examples}
\begin{itemize}[nosep]
  \item $x\mapsto\abs x$ and the ReLU $x\mapsto\max(0,x)$ are $1$-Lipschitz,
        although they are not differentiable at $0$.
  \item The sigmoid $\sigma(x) = 1/(1+e^{-x})$ has $\sigma' = \sigma(1-\sigma)\leq\frac14$,
        so it is $\frac14$-Lipschitz; $\tanh$ is $1$-Lipschitz.
  \item $x\mapsto x^2$ is not Lipschitz on $\R$ (its slope is unbounded), but it
        is $2a$-Lipschitz on $[-a,a]$.
  \item $x\mapsto\sqrt x$ on $[0,1]$ is continuous but not Lipschitz: near $0$
        the chord slopes $\sqrt h/h = 1/\sqrt h$ are unbounded.
  \item A linear map $x\mapsto Wx$ is Lipschitz with constant equal to its
        \emph{spectral norm} $\norm{W}_2 = \sigma_{\max}(W)$, the largest
        singular value, since $\norm{Wx_1 - Wx_2} = \norm{W(x_1-x_2)}$ and
        $\norm{W}_2$ is by definition $\max_{u\neq0}\norm{Wu}/\norm u$.
\end{itemize}
\end{example}

\section{Computing Lipschitz constants}\label{lipschitz-continuity:sec:computing}

\begin{proposition}[Lipschitz constant of a differentiable function]\label{lipschitz-continuity:prop:jacobian}
Let $f:\R^n\to\R^m$ be continuously differentiable with Jacobian $J_f$. Then
\[
  \mathrm{Lip}(f) = \sup_{x}\,\norm{J_f(x)}_2 .
\]
For $m=1$ this reads $\mathrm{Lip}(f) = \sup_x\norm{\nabla f(x)}$.
\end{proposition}

\begin{proof}
($\leq$) For $x_1,x_2$, let $\gamma(s) = x_1 + s(x_2-x_1)$. By the fundamental
theorem of calculus along the segment,
\[
  f(x_2) - f(x_1) = \int_0^1 J_f(\gamma(s))\,(x_2 - x_1)\,ds,
  \quad\text{so}\quad
  \norm{f(x_2)-f(x_1)} \leq \sup_x\norm{J_f(x)}_2\,\norm{x_2-x_1}.
\]
($\geq$) Fix $x$ and a unit vector $u$. For small $h>0$,
$\norm{f(x + hu) - f(x)}/h\leq\mathrm{Lip}(f)$, and the left side tends to
$\norm{J_f(x)u}$. Taking the supremum over $u$ gives
$\norm{J_f(x)}_2\leq\mathrm{Lip}(f)$.
\end{proof}

\begin{proposition}[Composition]\label{lipschitz-continuity:prop:composition}
$\mathrm{Lip}(f\circ g)\leq\mathrm{Lip}(f)\,\mathrm{Lip}(g)$.
\end{proposition}

\begin{proof}
$\norm{f(g(x)) - f(g(y))}\leq\mathrm{Lip}(f)\norm{g(x) - g(y)}\leq\mathrm{Lip}(f)\mathrm{Lip}(g)\norm{x-y}$.
\end{proof}

\begin{example}[A bound for a neural network]\label{lipschitz-continuity:ex:network}
A feed-forward network
$F(x) = W_K\,\varphi(W_{K-1}\cdots\varphi(W_1x))$ with a $1$-Lipschitz activation
$\varphi$ (ReLU, $\tanh$) applied coordinatewise is a composition of linear maps
and $1$-Lipschitz maps (biases do not change Lipschitz constants). By
\cref{lipschitz-continuity:prop:composition},
\[
  \mathrm{Lip}(F)\leq\prod_{k=1}^K\norm{W_k}_2 .
\]
The bound is cheap but can be very loose, because it ignores how the layers
interact. With $W_1 = \diag(1,0)$, $W_2 = \diag(0,1)$ and no activation, the
product of norms is $1$ while $W_2W_1 = 0$ is $0$-Lipschitz. Computing the exact
Lipschitz constant of a ReLU network is NP-hard.
\end{example}

\paragraph{Spectral normalization.} The spectral norm of a weight matrix is
estimated cheaply by \emph{power iteration}: repeat
$u\leftarrow W^\top Wu/\norm{W^\top Wu}$; then $\norm{Wu}\to\sigma_{\max}(W)$.
Dividing each weight matrix by this estimate during training (Miyato et al.,
2018) makes every layer $1$-Lipschitz, and hence, by the bound above, the whole
network.

\section{Robustness to adversarial perturbations}\label{lipschitz-continuity:sec:robustness}

Deep classifiers are vulnerable to \emph{adversarial examples}: inputs changed
by an imperceptible amount that receive a different label. A Lipschitz bound
turns ``imperceptible changes have small effects'' into a theorem.

\begin{proposition}[Certified radius]\label{lipschitz-continuity:prop:certified}
Let $F:\R^n\to\R^K$ output one score per class, be $L$-Lipschitz in the
Euclidean norm, and predict $c = \argmax_k F_k(x)$ with margin
\[
  \gamma(x) = F_c(x) - \max_{k\neq c}F_k(x) > 0 .
\]
Then the prediction is unchanged for every perturbation $\delta$ with
$\norm\delta < \gamma(x)/(\sqrt2\,L)$.
\end{proposition}

\begin{proof}
For $k\neq c$, the function $x\mapsto F_c(x) - F_k(x) = (e_c - e_k)^\top F(x)$
is a composition of $F$ with a linear map of norm $\norm{e_c - e_k} = \sqrt2$,
hence $\sqrt2L$-Lipschitz. Therefore
\[
  F_c(x+\delta) - F_k(x+\delta)\geq F_c(x) - F_k(x) - \sqrt2L\norm\delta
  \geq\gamma(x) - \sqrt2L\norm\delta > 0,
\]
so class $c$ still has the largest score.
\end{proof}

\begin{example}[A certified cat]
A $1$-Lipschitz classifier outputs class probabilities, and gives an image the
label ``cat'' with probability $0.99$; every other class has probability at
most $0.01$, so the margin is at least $0.98$. An adversary changes the image
by $\norm{\delta} = 0.1$. The output vector moves by at most $0.1$ in Euclidean
norm, so every coordinate moves by at most $0.1$: ``cat'' keeps a probability of
at least $0.89$ and every other class stays below $0.11$. The label cannot flip.
\Cref{lipschitz-continuity:prop:certified} gives the same conclusion for any
$\norm\delta < 0.98/\sqrt2\approx0.69$.
\end{example}

Training with \emph{Lipschitz regularization} (penalizing large gradients) or
with spectrally normalized layers keeps $L$ small and so enlarges the certified
radius; the price is usually some loss of accuracy on clean data.

\section{Smoothness and gradient descent}\label{lipschitz-continuity:sec:optimization}

In optimization the Lipschitz assumption is placed on the \emph{gradient} of the
loss.

\begin{definition}[Smoothness]
A differentiable function $\mathcal J:\R^p\to\R$ is \emph{$\beta$-smooth} if its
gradient is $\beta$-Lipschitz:
$\norm{\nabla\mathcal J(\theta_1) - \nabla\mathcal J(\theta_2)}\leq\beta\norm{\theta_1-\theta_2}$.
For twice-differentiable $\mathcal J$ this means that the Hessian has spectral
norm at most $\beta$ everywhere: the curvature is bounded.
\end{definition}

\begin{lemma}[Descent lemma]\label{lipschitz-continuity:lem:descent}
If $\mathcal J$ is $\beta$-smooth, then for all $\theta,\theta'$,
\[
  \mathcal J(\theta')\leq\mathcal J(\theta) + \nabla\mathcal J(\theta)^\top(\theta'-\theta) + \frac\beta2\norm{\theta'-\theta}^2 .
\]
In particular, one gradient step $\theta' = \theta - \eta\nabla\mathcal J(\theta)$
satisfies
\[
  \mathcal J(\theta')\leq\mathcal J(\theta) - \eta\Bigl(1 - \frac{\beta\eta}{2}\Bigr)\norm{\nabla\mathcal J(\theta)}^2,
\]
so the loss decreases whenever $0<\eta<2/\beta$ and $\nabla\mathcal J(\theta)\neq0$.
\end{lemma}

\begin{proof}
Let $v = \theta'-\theta$. By the fundamental theorem of calculus,
\[
  \mathcal J(\theta') - \mathcal J(\theta) - \nabla\mathcal J(\theta)^\top v
  = \int_0^1\bigl(\nabla\mathcal J(\theta + sv) - \nabla\mathcal J(\theta)\bigr)^\top v\,ds
  \leq\int_0^1\beta s\norm v^2\,ds = \frac\beta2\norm v^2,
\]
using Cauchy--Schwarz and the Lipschitz property. For the gradient step,
$v = -\eta\nabla\mathcal J(\theta)$, and the right-hand side becomes
$-\eta\norm{\nabla\mathcal J}^2 + \frac{\beta\eta^2}{2}\norm{\nabla\mathcal J}^2$.
\end{proof}

The lemma says that a smooth function lies below a parabola touching it at the
current point; gradient descent minimizes that parabola. The best guaranteed
decrease, $\frac1{2\beta}\norm{\nabla\mathcal J}^2$, is obtained for $\eta = 1/\beta$,
the classical choice of learning rate.

\begin{example}[The threshold $2/\beta$ is sharp]
For $\mathcal J(\theta) = \frac a2\theta^2$ with $a>0$, the gradient $a\theta$ is
$a$-Lipschitz, so $\beta = a$. Gradient descent gives
$\theta_{k+1} = \theta_k - \eta a\theta_k = (1 - \eta a)\theta_k$, which
converges to the minimizer $0$ iff $\abs{1-\eta a}<1$, i.e.\ $0<\eta<2/a$. At
$\eta = 2/a$ the iterates oscillate between $\pm\theta_0$ forever, and beyond it
they diverge. A learning rate that is too large ``overshoots'' because it
ignores the curvature.
\end{example}

\section{Other uses in machine learning}\label{lipschitz-continuity:sec:other}

\paragraph{Wasserstein GANs.} The Wasserstein-1 distance between distributions
$P$ and $Q$ has the dual form (Kantorovich--Rubinstein)
\[
  W_1(P,Q) = \sup_{\mathrm{Lip}(D)\leq1}\ \E_{X\sim P}[D(X)] - \E_{X\sim Q}[D(X)] .
\]
A Wasserstein GAN trains a \emph{critic} $D$ to approximate this supremum, so
the critic must be kept $1$-Lipschitz: by clipping weights (Arjovsky et al.,
2017), by penalizing $(\norm{\nabla D(\hat x)} - 1)^2$ at points $\hat x$
interpolating real and generated samples (WGAN-GP, Gulrajani et al., 2017), or
by spectral normalization. Without the constraint the supremum is infinite and
training is unstable.

\paragraph{Generalization.} Several generalization bounds for neural networks
depend on the network through margin-normalized products of layer norms, such
as $\prod_k\norm{W_k}_2$ (Bartlett, Foster and Telgarsky, 2017). Informally, a
function with a small Lipschitz constant cannot fit arbitrary labels on nearby
points, so it is less able to memorize noise. Such bounds explain trends more
than they predict test errors.

\paragraph{Reinforcement learning.} Assuming the value function and the policy
are Lipschitz in the state formalizes the idea that similar states have
similar values and call for similar actions, which is what allows
generalization across a continuous state space.

\paragraph{Metric learning.} If an embedding network $\phi$ is $L$-Lipschitz,
then $\norm{\phi(a) - \phi(b)}\leq L\norm{a-b}$: inputs that are close stay
close in the embedding, so similarity scores computed from embeddings vary
smoothly with the inputs.

\begin{table}[ht]
\centering
\caption{Roles of Lipschitz continuity.}
\label{lipschitz-continuity:tab:summary}
\begin{tabular}{@{}p{0.2\textwidth}p{0.24\textwidth}p{0.48\textwidth}@{}}
\toprule
Area & What is Lipschitz & Consequence \\
\midrule
robustness & the classifier $F$ & certified radius $\gamma/(\sqrt2L)$ \\
optimization & the gradient $\nabla\mathcal J$ & descent for $\eta<2/\beta$; $\eta = 1/\beta$ \\
WGAN & the critic $D$ & critic estimates the Wasserstein distance \\
generalization & the network & complexity measured by products of norms \\
RL, metric learning & value, policy, embedding & similar inputs give similar outputs \\
\bottomrule
\end{tabular}
\end{table}

\begin{chaptersummary}
\begin{itemize}[nosep]
  \item $f$ is $L$-Lipschitz if $\norm{f(x)-f(y)}\leq L\norm{x-y}$; for
        differentiable $f$, $\mathrm{Lip}(f) = \sup_x\norm{J_f(x)}_2$.
  \item Lipschitz constants multiply under composition; a network with
        $1$-Lipschitz activations satisfies
        $\mathrm{Lip}(F)\leq\prod_k\norm{W_k}_2$, a cheap but loose bound
        that spectral normalization controls.
  \item A margin $\gamma$ certifies robustness to perturbations of norm below
        $\gamma/(\sqrt2L)$.
  \item A $\beta$-Lipschitz gradient gives the descent lemma: gradient steps
        decrease the loss for $\eta<2/\beta$, and this threshold is sharp.
  \item The Kantorovich--Rubinstein duality is why WGAN critics must be
        $1$-Lipschitz.
\end{itemize}
\end{chaptersummary}

\section*{Exercises}
\addcontentsline{toc}{section}{Exercises}

\begin{exercise}\label{lipschitz-continuity:exr:1}
Show that $\log(1 + e^{x})$ (the softplus) is $1$-Lipschitz and that its
derivative is $\frac14$-Lipschitz. Deduce a safe learning rate for the
one-dimensional logistic loss $\theta\mapsto\log(1 + e^{-y\theta x})$ with
$\abs x\leq 1$, $y\in\{\pm1\}$.
\solutionlink{lipschitz-continuity:exr:1}
\end{exercise}

\begin{exercise}\label{lipschitz-continuity:exr:2}
Compute the Lipschitz constant of $x\mapsto Wx$ for
$W = \begin{psmallmatrix}3 & 0\\ 4 & 5\end{psmallmatrix}$, and compare with
the Frobenius norm $\norm W_F$. Which is the larger, and why is that always the
case?
\solutionlink{lipschitz-continuity:exr:2}
\end{exercise}

\begin{exercise}\label{lipschitz-continuity:exr:3}
Let $\mathcal J(\theta) = \frac12\theta^\top A\theta$ with $A$ symmetric
positive definite and eigenvalues $0<\lambda_{\min}\leq\lambda_{\max}$. Show that
$\mathcal J$ is $\lambda_{\max}$-smooth and that gradient descent converges iff
$0<\eta<2/\lambda_{\max}$. Which eigen-direction converges slowest with
$\eta = 1/\lambda_{\max}$?
\solutionlink{lipschitz-continuity:exr:3}
\end{exercise}
\end{notechapter}
